\section*{Chapitre \ref{ch:s2:credit-relative-value} --- Valeur relative du crédit}

\begin{solution}{exo:s2:credit-relative-value:1}
Par unité de notionnel, $15 - 0.9 \times 30 = -12$ points de base par an~; sur un capital de 10~\% cela représente $-1.2~\%$ par an.
\end{solution}

\begin{solution}{exo:s2:credit-relative-value:2}
$4.5 \times 1.35~\% = 6.1~\%$ du notionnel, soit $60.8~\%$ du capital avec un levier de dix.
\end{solution}

\begin{solution}{exo:s2:credit-relative-value:3}
$1/0.10 = 10$ fois.
\end{solution}

\begin{solution}{exo:s2:credit-relative-value:4}
En cas de défaut, la protection verse le pair moins le taux de recouvrement et l'obligation verse le recouvrement~: la position est intégralement compensée. Avant cela, l'obligation et le CDS sont valorisés
à leurs spreads, et lorsque la base s'élargit, l'obligation perd davantage que la protection ne gagne~; avec l'effet de levier, les valorisations peuvent forcer une vente.
\end{solution}

\begin{solution}{exo:s2:credit-relative-value:5}
Des frais de prêt élevés signifient que de nombreux investisseurs sont vendeurs à découvert de l'obligation~; les arbitragistes hésitent à détenir une obligation contre laquelle d'autres parient, même couverte, de sorte que la
base doit être plus négative pour les attirer.
\end{solution}

\begin{solution}{exo:s2:credit-relative-value:6}
Un défaut entre les deux échéances~: un flattener vend une protection à courte échéance et en achète une à longue échéance~; un défaut survenant avant l'échéance du contrat court
verse les deux jambes, qui se compensent à peu près, tandis que les notionnels, dimensionnés par la duration de spread, sont inégaux et laissent une exposition nette au jump-to-default.
\end{solution}

\begin{solution}{exo:s2:credit-relative-value:7}
Sans coût de financement, vendre après une perte cumulée de 50~\% du capital fige un résultat de $-50.5~\%$ à l'année 5.24, quelques jours avant le point bas de la base~; conserver la position jusqu'au terme
rapportait $+24.4~\%$. La règle protège le portefeuille de la ruine dans une crise plus grave, mais lui fait renoncer au rebond de celle-ci~; c'est une décision de levier prise après
coup.
\end{solution}

\begin{solution}{exo:s2:credit-relative-value:8}
Détenir jusqu'à l'échéance suppose un financement jusqu'à l'échéance. L'obligation est financée en repo à court terme, avec une décote susceptible de monter, et les appels de marge sur les valorisations
forcent des ventes avant l'échéance~; l'arbitrage n'est réel que pour le capital qui ne peut être retiré.
\end{solution}

\begin{solution}{pb:s2:credit-relative-value:1}
\begin{enumerate}
\item Spread du CDS moins spread de l'obligation~; achat de l'obligation avec protection lorsque le spread de l'obligation est plus élevé.
\item L'obligation couverte est sans risque~; les frictions (financement, liquidité, contrepartie, collatéral) empêchent de l'arbitrer.
\item Les écarts les plus importants concernent les obligations les plus frictionnelles, et sont plus négatifs lorsque les frais de prêt sont élevés.
\item Une prime de 15 points de base, portée à 150 en l'année 5 pendant un trimestre, puis ramenée à son niveau initial en un an.
\item Indice contre CDS des membres lorsque le skew est large~; 55 points de base par an, ratio de Sharpe 0.97.
\item Deux échéances, à duration égale~; 49 points de base par an, ratio de Sharpe 0.75.
\item Coûts, événements de crédit des entités de référence, jump-to-default entre les échéances.
\item Il faut quarante transactions sur noms individuels.
\item $+1.6~\%$, $-1.1~\%$ et $-3.8~\%$ par an pour un financement à 0, 30 et 60 points de base.
\item Environ $-60~\%$ du capital jusqu'au point bas~; $+76~\%$, $+63~\%$, $+50~\%$ ensuite.
\item Effet de levier retiré par les prime brokers défaillants~; arbitragistes contraints de vendre.
\item $+69.8~\%$ sans coût de financement, $+67.1~\%$ à 30 points de base, sur un an.
\item \textbf{Résultat nommé.} En temps normal, le \omterm{def:s2:credit-relative-value:negbasis}{trade de base négative} rapporte 1.6~\% par an sur le capital à coût de financement nul et perd de l'argent à 30 points de
base~; lorsque les coûts de financement bondissent et que la base s'élargit de $-15$ à $-151$, il perd environ 60~\% de son capital avec un levier de dix, le tout par la valorisation.
\item Une base de $-15$ couvre peu de financement.
\item Capital qui n'était pas porté à la limite de levier avant la crise.
\item Pour une crise qui élargit la base, au levier que le prêteur accepte encore alors.
\item Financement et décotes tels qu'ils étaient, ventes forcées, la crise de l'échantillon.
\item Crédit cross-currency.
\item Les trois sont des arbitrages financés dont les pertes viennent du financement, et non de la couverture.
\item Le coût du bilan qui détient l'obligation.
\end{enumerate}
\end{solution}

\begin{solution}{iq:s2:credit-relative-value:1}
Le spread du CDS moins le spread de l'obligation pour le même émetteur~; négatif signifie que l'obligation rapporte plus que la prime du CDS, de sorte qu'acheter l'obligation avec protection
procure un spread positif, qui persiste parce que le financement et la détention de l'obligation sont coûteux.
\end{solution}

\begin{solution}{iq:s2:credit-relative-value:2}
Acheter l'obligation, la financer en repo, acheter la protection~; percevoir les coupons de l'obligation, payer le repo et la prime du CDS~; en cas d'événement de crédit, livrer l'obligation (ou régler en
espèces) et recevoir le pair~: la perte sur l'obligation est couverte et la position est clôturée.
\end{solution}

\begin{solution}{iq:s2:credit-relative-value:3}
Un élargissement de la base avec des décotes et des taux de repo plus élevés, une défaillance de contrepartie sur le CDS, un squeeze à la livraison et la perte du
financement du portefeuille~; à un levier de dix, le scénario de crise du chapitre.
\end{solution}

\begin{solution}{iq:s2:credit-relative-value:4}
Le spread intrinsèque est la moyenne des spreads des membres, pondérée de sorte que la jambe de prime de l'indice corresponde à celle du panier (à peu près la moyenne pondérée par la duration)~;
l'indice s'en écarte parce qu'il est plus liquide, qu'il se négocie en premier et qu'il sert à la couverture.
\end{solution}

\begin{solution}{iq:s2:credit-relative-value:5}
Entités de référence et obligations, séries et versions d'indices, conventions de recouvrement, coupons et dates de roll~; à chaque roll, l'indice on-the-run change de
membres et d'échéance, et les positions doivent être rattachées à la bonne série.
\end{solution}

\begin{solution}{iq:s2:credit-relative-value:6}
Avant défaut, la position reçoit le spread de l'obligation au-dessus du taux sans risque et verse la prime du CDS~; en cas de défaut, elle reçoit le pair de la protection
et rien de plus de l'obligation que le recouvrement qu'elle livre. Le financement, le collatéral, l'option de moins cher à livrer (CTD) et le risque de contrepartie
rompent l'équivalence.
\end{solution}
